\section[第7章\q 第六课\q 学会不为钱工作]{第7章\\第六课\q 学会不为钱工作}

1995年，我接受了新加坡一家报纸的采访。一位年轻的女记者准时赴约，采访立即开始了。我们坐在一家豪华酒店的大厅里，喝着咖啡，谈论我此次行程的目的。我和畅销书作家金克拉一起接受采访，他谈的是动机问题，而我谈的是“富人的秘密”。

“有一天，我想成为像你这样的畅销书作家。”女记者说。我曾经读过她在报上发表的一些文章，而且印象深刻，她的文章风格犀利且有条理，深受读者的欢迎。

“你的文章风格很好，”我回答，“那么，是什么妨碍了你实现梦想？”

“我的写作之路似乎前途渺茫，”她平静地说，“人们都说我的小说非常好，但是仅此而已。因此，我依然继续在报社工作，至少，这能挣钱支付账单。你能给我些建议吗？”

“好的，”我高兴地说，“我在新加坡的一位朋友办了一所学校，培训销售技能。他为新加坡的许多大公司讲授营销课程，我想如果你去听听他的课，或许会对你的职业生涯有帮助。”

她有点不快，“你是说我应该去学销售吗？”

我点点头。

“你是当真的吗？”

我又点点头。“有什么不妥吗？”我问道。但她似乎被激怒了。我现在有点后悔自己所说的话。本来我是想帮忙，现在却得为自己的建议辩解。

“我拥有英语文学硕士学位，我干吗要去学做推销员？我是一个专业人士，即使我需要接受再教育也是为了进行专业上的深造，绝不是为了去当什么推销员。我讨厌那些推销员，他们眼里只有钱。所以，请告诉我为什么我非得去学销售？”她边说边用力地抓起她的提包，于是采访草草收场了。

在咖啡桌旁放着我以前写的一本畅销书。我拿起书，见到她粘在封面上的一张便条。“你看到这个了吗？”我指着她记的便条。

她低头去看自己的便条，“什么？”她困惑地说。

我又指了指她的便条，在便条上写着：“罗伯特·T.清崎，畅销书作家。”

“上面写的是畅销书作家，而不是最好的作家。”

她突然睁大了眼睛。

“我只是一个平庸的作家，而你则是一位优秀的作家。我去了销售学校，而你得了硕士学位。如果你把这两方面结合起来，你就既是‘畅销书作家’，又是‘最好的作家’。”

她的眼里燃起了怒火，“我永远不会委屈自己去学什么销售，像你这样的人也不应该写作。我是受过专业训练的作家，而你以前不过是一个推销员，这不公平。”

她收好其他的便条，然后匆匆穿过巨大的玻璃门，消失在新加坡潮湿的清晨里。

好在第二天早上她给了我一个公平、良好的访谈记录。

全世界到处都有聪明、才华横溢、受过良好教育以及极具天赋的人。我们每天都会碰到这样的人，他们就在我们的身边。

几天前，我的汽车出了点问题。我把它开进维修厂，一个年轻的机械工不到几分钟就修好了。他只要听听发动机的声音就能确定毛病出在哪里，这使我感到非常惊讶。

然而遗憾的是，只有天赋是不够的。

我常常吃惊为什么有些人才华过人却只挣到很低的工资。我听说只有不到5％的美国人年收入超过10万美元。我曾遇见聪明且受过高等教育的人每年收入不到2万美元。一位药品贸易的商务顾问曾经告诉我，有许多医生、牙医和按摩师在财务上困难重重。以前我总是认为他们一毕业，美元就会滚滚而来。这位商务顾问还跟我说了一句话：“他们只有一项技能，所以挣不到大钱。”

这句话的意思是说，大部分人需要学习和掌握不止一项技能，只有这样他们的收入才能获得显著增长。以前我提到过，财商是会计、投资、市场和法律等各方面知识和能力的综合。将上述4种技能结合起来，以钱生钱就会容易得多。当涉及钱的时候，只有一项技能的人不得不努力工作。

有关综合技能的典型例子就是那位为报纸撰稿的年轻女记者。如果她能努力学习销售的技能和有关市场的知识，她的收入就会显著增加。

如果我是她，我一定会去学习一些有关广告文案和销售方面的课程。然后，我会去一家广告公司找一份工作。即使这样做会使收入减少，但我能学到在成功的广告中运用的“用几秒钟交流”的技巧。我还会花时间去学习公共关系这一重要的技能，以便通过免费的公关来赚取数百万美元。然后，利用晚上和周末来创作。如果她能够做到我上面说的事，她写的书就会畅销，并且，会在很短的时间内，成为一位“畅销书作家”。

当我第一次带着我写的书《如果你想生活得富裕幸福，要不要去学校？》去见一位出版商时，他建议我将书改名为《经济学教育》。我告诉出版商，如果用这个书名，我只能卖出两本书：一本给我的家人，另一本给我最好的朋友。可问题是，即便是他们也希望免费得到这本书。

选择《如果你想生活得富裕幸福，要不要去学校？》这一“可恨的”书名，是因为我们知道它会受到大众的欢迎。我赞成教育，但认为应该进行教育改革，否则，我为什么一直在呼吁改革陈旧的教育体制呢？我之所以选择这样一个书名，使我有机会在更多的电视和电台节目中露面，因为我愿意引起争议。许多人可能认为我没什么深度，这本书却一版再版。

1969年，我从美国商船学院毕业了。我受过良好教育的爸爸十分高兴，因为加州标准石油公司录用了我，让我在油轮上工作。我是一个三副，比起我的同班同学，我的工资很低，但作为我离开大学之后的第一份工作，也还算不错。我的起薪是一年4.2万美元，包括加班费。而且我一年只需工作7个月，余下的5个月是假期。如果我愿意的话，还可以用那5个月的假期随一家附属船舶运输公司的船只到越南去，这样年收入能翻一番。

尽管前途光明，但我还是在6个月后辞职了，加入海军陆战队去学习飞行。对此我受过良好教育的爸爸非常失望，富爸爸则祝贺我作出了这样的决定。

在学校和工作单位，最普遍的观点就是“专业化”。也就是说，为了挣更多的钱或者能升职，你需要更加专业化。这就是医生们要早早寻求某种专长——如骨科或儿科——的原因。对于会计师、建筑师、律师、飞行员及其他很多专业人士也是这样。

我那受到良好教育的爸爸也信奉同样的教条，因此，当他最终获得博士学位时非常激动。不过他也承认，社会越来越少奖励那些学得多的人了。

富爸爸则鼓励我去做恰好相反的事情。“对许多知识你只需知道一点就够了”，这是他的建议。所以，这些年来我曾在他的公司的不同部门工作过。有一段时间我在他的会计部工作，虽然我从来不想成为会计，但他仍希望我借助这种渗透作用学习一些会计方面的常识。富爸爸知道我会明白那些“行话”，而且懂得哪些东西重要，哪些东西不重要。

我也曾做过公共餐厅服务员、建筑工人、推销员、仓库保管员和市场营销人员。富爸爸一直在培养我和迈克，坚持让我们列席他与银行经理、律师、会计师和经纪人的会议，希望我们对他的商业帝国的每一部分都能有所了解。

当我辞掉在标准石油公司收入丰厚的工作后，受过良好教育的爸爸和我进行了推心置腹的交谈。他非常吃惊，也不理解我为什么要辞去这样一份工作：收入高，福利待遇好，休假长，还有升职的机会。有一天晚上他问我：“你为什么要辞职呢？”我没法向他解释清楚，我的逻辑与他的不一样。最大的问题就在于此，我的逻辑和富爸爸的一致。

对于受过良好教育的爸爸来说，工作的稳定就是一切；而对于富爸爸来说，不断学习才是一切。

受过良好教育的爸爸认为我去学校学习就是要做一名船员，而富爸爸则知道我去学校是为了学习国际贸易。因此，在我做学生时就跑过货运，驾驶过去远东及南太平洋的大型运输船、油轮和客轮。富爸爸强调我应该乘船去太平洋而不是去欧洲，因为他认为“新兴国家”是在亚洲而不是欧洲。当我的大多数同班同学，包括迈克，在搞联谊的时候，我却在日本、泰国、新加坡、越南、韩国、菲律宾以及中国台湾和中国香港等地学习贸易、人际关系、商业类型和当地文化。我也参加晚会，但不去任何大学的联谊会，我迅速地成熟起来了。

受过良好教育的爸爸无法理解我为什么决定辞职加入海军陆战队。

我告诉他我想要学习飞行，但实际上我是想学会指挥军队。富爸爸曾给我解释，经营一家公司最困难的就是对人员进行管理。他在军队待过3年，而受过良好教育的爸爸则免服兵役。富爸爸告诉我学会在危险形势下领导下属的重要性。“领导才能是你下一步迫切需要学习的，”他说，“如果你不是一个好的领导者，你就会被背后的冷箭射中，就像他们在商业活动中做的一样。”

1973年从越南回国后，我离开了军队，尽管我仍然热爱飞行。我在施乐公司找到了一份工作，我是有目的的，不过不是为了物质利益。我是一个腼腆的人，对我而言销售是世界上最令人害怕的课程，而施乐公司的营销培训项目是美国最好的之一。

富爸爸为我感到自豪，而受到良好教育的爸爸则为我感到羞愧。作为知识分子，他认为推销员低人一等。我在施乐公司工作了4年，直到我不再为可能吃闭门羹而发怵。当我稳居销售业绩榜前5名时，我再次辞职，又放弃了一份不错的职业和一家优秀的公司。

1977年，我创建了自己的第一家公司。富爸爸鼓励我和迈克去管理公司，现在我就得学着创建并管理它们了。我的公司的第一种产品是尼龙和维可牢\footnote{维可牢经常用于布制品上，由一条表面有细小钩子的尼龙条与表面有毛圈的对应的尼龙条粘合面构成。}搭扣钱包，它们是在远东生产的，然后装船运到纽约的仓库里，仓库离我去上学的地方很近。我的常规教育已经完成，现在是我单飞的时候了。如果失败了，我就会破产。富爸爸认为如果要破产的话，一定要在30岁以前，他的建议是“这样你还有时间东山再起”。就在我30岁生日的前夜，我的第一批货物驶离韩国前往纽约。

直到今天，我仍然在做国际贸易，就像富爸爸鼓励我的那样，我一直在寻找新兴国家的商机。现在我的投资公司在南美、亚洲、挪威和俄罗斯等地都拥有投资项目。

正如一句格言所说：“工作(job)一词就是‘比破产强一点’(just over broke)。”然而不幸的是，这句话确实适用于千百万人，因为学校没有把财商看做是一种智慧，大部分工薪阶层都量入为出，他们挣钱，然后支付账单。

还有另外一种可怕的管理理论是这样说的：“工人付出最大努力以免于被解雇，而雇主提供最低工资以防止工人辞职。”如果你看一看大部分公司的工资支付额度，就会明白这一说法确实在某种程度上道出了真相。

最终的结果是大部分人从不敢越雷池一步，他们按照别人教他们的那样去做：找一份稳定的工作。大部分人是为短期的工资和福利工作的，但从长期来看这种做法常常是具有灾难性的。

相反，我劝告年轻人在找工作时要看能从中学到什么，而不是只看能挣多少钱。在选择某种职业或陷入“老鼠赛跑”的陷阱之前，要仔细看看脚下的路，弄清楚自己到底想获得什么技能。

一旦人们陷入为支付账单而整天疲于奔命的陷阱，就和那些蹬着小铁笼不停转圈的小老鼠一样了。它们的小毛腿蹬得飞快，小铁笼也转得飞快，可到了第二天早上，它们发现自己依然被困在笼子里，就像你被你的工作困住一样。

在巨星汤姆·克鲁斯主演的电影《甜心先生》中，有许多非常好的台词。可能最容易记住的就是那句“把钱给我看看”。但我认为还有一句台词简直就可以被称为真理。他发生在汤姆·克鲁斯离开公司的那一幕，他刚被炒了鱿鱼，于是就问公司的人：“谁愿意和我一起走？”顿时鸦雀无声，仿佛连空气都凝固了。只有一位女士站出来说：“我愿意……可是3个月后我就能升职了。”

这句话可能是整部电影里最实在的一句台词，道出了那些总是为生计而忙碌工作的人们的心声。我知道，受到良好教育的爸爸每年都期望加薪，但每年都十分失望。于是他不得不回学校去获得更高的学历，以得到另一次加薪的机会，但是他只能又一次失望。

我经常向人们提一个问题：“你终日忙碌的目的是什么？”就像那只小老鼠一样，我想知道人们是否会想一想，这样辛苦地工作到头来究竟是为了什么？未来的日子又会怎样呢？

美国退休者协会前会长西里尔·布里克弗利克的报告说，个人退休金管理正处于一种混乱的状态。首先，在今天有50％的劳动力没有退休金，而在另外的50％的人中，有75％～80％的人的退休金不能足额发放，他们每月只能领到55美元、150美元或是300美元。

克莱格·卡佩尔在《退休的迷思》一书中写道：我采访过一家重要的全国性退休金咨询公司，并与一位专门为企业高管制定退休计划的经理进行了一次谈话。当我问她职业人士对养老金有什么期待时，她自信的笑了笑说：“万灵丹。”

“什么是万灵丹？”我问道。

她耸耸肩，说：“如果婴儿潮一代发现，在年老的时候并没有足够的钱来维生，他们会陷入绝望。”卡佩尔接着分析了原来的固定福利计划和后来更加不可靠的401(k)计划之间的区别。对于今天仍在工作的大部分人来说，这可不是一幅美妙的图景，而这仅仅是指退休金，如果再加上医疗和私人养老院的费用，这幅图景会更加可怕。在这本书中，卡佩尔指出每年的私人养老院的费用平均高达3万～12.5万美元。1995年，当他去当地的一家普普通通的私人养老院时，发现价格竟达到每年8.8万美元。

在一些拥有社会医疗保障的国家，许多医院不得不艰难地作出抉择，例如：到底先救谁。他们完全是根据这些病人有多少钱、年纪多大而作决定的。如果病人年纪大了，他们就会把医疗服务提供给更年轻的人，而那些又老又穷的病人只好排在最后。因此，就像富人能得到更好的教育一样，富人也能使自己活得更长，而穷人只会死得早一些。

所以我想知道，是否工薪阶层只有在梦想着未来或者等到下次领工资的时候，才会对自己的遭遇产生疑问呢？

当我面对那些想挣更多钱的成年人演讲时，我总是建议他们要有长远的眼光。我承认为了金钱和生活安稳而工作是很重要，但我仍主张要再找一份工作，以便从中学到另一种技能。我常常提议，如果想学习销售技能，最好进一家网络营销公司，也被称为多级营销公司。这类公司多半能够提供良好的培训项目，帮助人们克服因失败造成的沮丧和恐惧心理，这种心理往往是导致人们不成功的主要原因。从长远来看，教育比金钱更有价值。

当我提出这些建议时，我常常听到这样的回答：“这太麻烦了”，或是“我只想做我感兴趣的事”。

对于“太麻烦了”的说法，我反问：“那么，你宁可辛苦一生，把挣来的50％的收入交给政府？”对于另一种说法“我只想做我感兴趣的事”，我的回答是：“我对去健身房不感兴趣，但还是要去，因为我想身体更好，活得更长。”

遗憾的是有一些老话仍然不过时，像“你无法教老狗学会新把戏”，除非一个人习惯于变化，否则改变对他来说将是十分困难的。

但是，你们中间仍有些人对于“工作是为了学习新东西”这种观点有疑虑，我想说一句话鼓励你们：生活就像去健身房，最痛苦的事情是作出锻炼身体的决定，一旦你过了这一关，以后的事情就好办了。有很多次，我害怕去健身房，但是只要我去了并开始运动，就会感到非常愉快。健身之后我总是很高兴，因为我说服了自己坚持了下来。

如果你坚持不愿意学习新东西，而只是想在自己的领域里成为专家，那么你要确信你服务的公司是有工会的，因为工会会保护专业人士。

我受过良好教育的爸爸在失去了政府的“恩宠”之后，成为了夏威夷教师工会的负责人。他告诉我这是他所做过的最难的工作。对我个人来说，我不倾向于劳资任何一方，因为我能理解双方有各自的需要和利益。如果你按学校所教育的那样去做，成为一位专业人士，那么最好去寻求工会的保护。例如，如果我继续我的飞行生涯，我就会进一家有强有力的飞行员工会的公司。为什么？因为我将终生只在这一行里学到一种有价值的技能，如果我被这一行业抛弃，我的技能对其他行业便毫无用处。一位拥有10万小时驾驶大型运输机记录的高级飞行员，每年能挣15万美元，可一旦下岗，就很难找到一个收入相当的在学校教书的工作了。技能不一定能从一个行业转到另一个行业，在航空业被看重的飞行技能在教育系统并不受重视。

甚至对于今天的医生来说也是如此。随着医学的变化，许多医药专家需要加入“健康维护组织”这样的医疗机构，教师也一定要成为工会的会员。在如今的美国，教师工会是所有工会中最大、最富有的一个。全国教育协会拥有巨大的政治影响力。教师们需要工会的保护，因为他们技能的价值也只限于教育系统。因此法则就是：如果你是高度专业化的人士，就加入工会。这才是明智之举。

当我问自己班上的学生“你们中间有多少人能做出比麦当劳更好的汉堡包”时，几乎所有的学生都举起了手。我接着问：“如果你们中大部分人都能做出比麦当劳更好的汉堡包，那为什么麦当劳比你们更赚钱？”

答案是显而易见的：麦当劳拥有一套出色的商业体系。许多才华横溢的人之所以贫穷，就是因为他们只是专心于做好产品，而对商业体系却知之甚少。

我有一位夏威夷的朋友是杰出的艺术家，他也挣了很多钱。一天，他母亲的律师打电话告诉他，他母亲给他留下了3.5万美元的遗产，这是他母亲的房产在扣除律师费和税收后的余额。不久后，他发现了一个可以发展事业的“机会”。为此，他需要用这笔钱的一部分来做广告。两个月后，他的第一个四色整版广告登在一份昂贵的杂志上，这份杂志的读者主要是富人。然而广告刊登了3个月后，却没有收到任何效果，他所继承的遗产也花光了。现在他想以误导为由起诉那家杂志。

这是一个有关只懂怎么做好汉堡包，却不懂如何将汉堡包卖出去的典型例子。当我问他从这件事学到了什么时，他只是回答“广告商都是骗子”。于是我问他愿不愿意学习销售和一门直销课程，他回答：“我没时间，也不愿意浪费钱。”

世界上到处都是有才华的穷人。在很多情况下，他们贫穷、财务困难或者只能挣到低于他们应得的薪水，不是因为他们已知的东西而是因为他们未知的东西。他们只将注意力集中在提高做汉堡包的技能上，却不注意提高销售和配送汉堡包的技能。也许麦当劳不能做出最好的汉堡包，但他们能在做出一般水平的汉堡包的前提下，做到最好的销售和配送工作。

穷爸爸希望我学有所长，这是在他看来能够获得更高收入的途径。

即使是在夏威夷州长通知他不能继续在政府工作时，我受到良好教育的爸爸仍然鼓励我去拥有某项专长。后来，受到良好教育的爸爸接手了教师工会的工作，为保护那些高级专业人才和受到良好教育的人士的利益而努力。我们经常为此争论不休，但我知道，他从不认为过分专业化恰恰是导致这些人需要工会保护的原因。他不能理解，为何越是专业化，就越是掉入陷阱而无法自拔。

富爸爸建议我和迈克去“培养”自己。许多企业也是这么做的，他们在商学院里挑选一个年轻聪明的学生，并开始“培养”他，希望有朝一日他能领导这家公司。因此，这些聪明的年轻人并不去专项钻研某一个部门的业务，而是从一个部门跳到另一个部门，从而学到整个企业各个系统的知识。富人们也常常这样“培养”自己的或是别人的孩子，通过这种方法，孩子们能对如何经营一家企业有一个整体的认识，并可以了解不同部门的相互关系。

对于经历过第二次世界大战的那一代人来说，从一家公司跳槽到另一家公司是一件坏事，而今天人们却认为这是明智之举。既然人们从一家公司跳到另一家公司，不是为了寻求更深入的专业知识，那为什么不借此机会多学习而别光想多挣钱呢？尽管从短期来看，你的薪水可能会减少；但从长远来看，你将从中获得巨大的收益。

成功所必需的管理技能包括：

{\kaishu
1. 对现金流的管理。

2. 对系统（包括你本人、时间及家庭）的管理。

3. 对人员的管理。}

最重要的专门技能是销售和对市场营销的理解。销售技能是个人成功的基本技能，它涉及与其他人的交往，包括与顾客、雇员、老板、配偶和孩子。而沟通能力，如书面表达、口头表达及谈判能力等对一个人的成功来说更是至关重要。我就是通过学习各种课程、听教学磁带等来扩展知识并不断提高自己的这一技能的。

正如我之前提到的那样，我受过良好教育的爸爸工作越努力，就越具有竞争力，但同时他也更深地陷入专业特长的陷阱之中。虽然他的工资增长了，可他的选择机会却少了。直到失去了在政府中的工作，他才发现自己在职业选择上是多么被动。这就好比职业运动员因为突然受伤或是年龄太大而无法继续参加比赛一样，他们会失去曾经拥有的高收入工作，而有限的技能又使他们无法另辟蹊径。我想，这就是为什么从那时起我爸爸变得如此支持工会了，因为他意识到工会能使他受益。

富爸爸鼓励我和迈克多去涉猎一些东西。他鼓励我们去和比我们更精明的人一起工作，并把他们组成一个团队。现在把这种做法称为专家组合。

今天，我看到一些曾当过老师的人现在每年能挣数十万美元，他们挣这么多钱是因为他们不仅拥有教育方面的专业技能，还拥有其他方面的特长。他们既能教书，也能做销售和市场营销。我还不知道有比销售和市场营销更重要的技能，但要掌握它们对大部分人来说是很困难的，这主要是因为他们害怕被拒绝。所以，你在处理人际关系、商务谈判和被拒绝时的恐惧心理等方面做得越好，生活就会越轻松。就像我对那位想成为“畅销书作家”的女记者所建议的一样，我今天也给所有人这个建议。精通专业技能既是优势也是弱点。我有许多朋友，他们非常有天赋，但不善于与其他人交流，结果他们的收入少得可怜。我建议他们用一年时间来学销售，即使挣不到什么钱，可他们处理人际关系的能力会大大提高，这种能力是无价的。

除了成为优秀的学习者、销售员和市场营销人员外，我们还需要成为好老师、好学生。要想真正富有，我们既要不吝付出也要学会索取。

对于那些被财务或职业问题所困的人来说，他们常常既不会付出，也无力索取。我知道许多人穷是因为他们既不是好老师也不是好学生。

我的两个爸爸都是很慷慨的人，他们都把付出放在第一位。向别人传授经验是他们付出的途径之一，他们付出的越多，得到的也就越多。

但他们有一个明显的区别，就是对金钱的付出。我的富爸爸会给别人许多钱，他把钱捐给教堂、慈善机构以及他的基金会。他知道要想得到金钱，就必须先付出金钱。付出金钱是那些非常富有的家庭保持富有的秘诀，也是例如洛克菲勒基金会、福特基金会这样的机构存在的原因。建立这些机构是为了增加财富，从而让其他人能够永远受益。

我那受过良好教育的爸爸总是说：“当我有多余的钱时，就会把它捐出来。”可问题是他从来就没有多余的钱。因此他更加努力地工作以增加收入，却没有注意到一条最重要的金钱法则——给予，然后获得。

相反，他信奉的是“得到后再付出”。

总之，我同时受到两个爸爸的影响。一方面我是资本主义的坚定信奉者，喜欢以钱生钱的游戏；另一方面我又是一个怀有社会责任感的老师，深切地关注日益加大的贫富差距。我个人认为，陈旧的教育体系应该对这一差距的加大负有重要责任。
